% ===================================================================
%  TABEL Nr. 5 — Die huis en sy bou.
%  Meme regime que les precedents.
%
%  CE TABLEAU EST UN DIALOGUE D'UN BOUT A L'AUTRE : trente-six
%  attributions de parole en « \textsc{...}. --- », comme dans l'ido.
%
%  LES NOMS PROPRES DE L'IDO SONT DES MOTS : « siorulo Blanko » et
%  « siorulo Rufo » sont M. Blanc et M. Roux. On les traduit — meneer
%  Wit, meneer Rooi — comme les prenoms.
%
%  ET C'EST ICI QUE LA REGLE DU « sch » TRAVAILLE LE PLUS. Le
%  neerlandais dit schoorsteen, schuur, schaaf, schilder, schroef ;
%  l'afrikaans skoorsteen, skuur, skaaf, skilder, skroef. Sept
%  occurrences dans ce seul tableau, toutes au metier du batiment —
%  c'est le vocabulaire ou les deux langues se ressemblent le plus, et
%  donc celui ou la lettre decide.
%
%  LE « z » INITIAL AUSSI : zolder devient solder, zink reste zink
%  (emprunt), zaag devient saag.
% ===================================================================

%%K t05-apar-1 apar
\VUsaut{6.0mm}
\VUtitre{15.0pt}{60}{TABEL Nr. 5}
\VUsaut{1.4mm}
\VUtitre{11.4pt}{0}{Die huis en sy bou.}
\VUsaut{2.2mm}

%%K t05-tit-1 sub
\VUpk{10.2pt}{40}{'n Gelukkige ontmoeting.}
\VUsaut{3.0mm}

%%K t05-01-1 p
\textsc{Jan}. --- Oom, ek wil graag weet hoe 'n mens 'n \VUgras{huis} bou.

%%K t05-01-2 p
\textsc{Die oom} (80). --- Dit is baie maklik, my vriend, kom saam met my, ek sal jou
die een wys wat meneer Bernardus (79) laat bou en waarin ek sal woon.

%%K t05-01-3 p
\textsc{Jan}. --- Wie het dan die \VUgras{plan} \textsuperscript{(76)} van hierdie huis
geteken?

%%K t05-01-4 p
\textsc{Die oom}. --- Die \VUgras{argitek} \textsuperscript{(75)}; hy het 'n baie
duidelike tekening voorgelê wat goedgekeur is, en die \VUgras{aannemer}
\textsuperscript{(77)} het die werk begin.

%%K t05-01-5 p
\textsc{Jan}. --- En wie is die aannemer?

%%K t05-01-6 p
\textsc{Die oom}. --- Dit is meneer Wit, die vader van jou vriend Jakobus. Hy lei die
werkers saam met meneer Rooi, die argitek.

%%K t05-01-7 p
Wel! Hier kom meneer Wit net betyds met sy \VUgras{liniaal} \textsuperscript{(78)}, en
meneer Rooi met sy plan.

%%K t05-01-8 p
Goeiedag, menere. Hoe gaan dit met u?

%%K t05-01-9 p
\textsc{Meneer Rooi}. --- Baie goed, dankie. Goeiedag, Jan, hoe het hierdie kind
gegroei!

%%K t05-01-10 p
\textsc{Die oom}. --- Ja, waarlik, hy is 'n klein man. Hy is baie ernstig, 'n mens moet
hom oor alles inlig wat hy sien. Vandag wil hy weet hoe 'n mens 'n huis bou.

%%K t05-tit-2 sub
\VUpk{10.2pt}{40}{Die werkers.}
\VUsaut{3.0mm}

%%K t05-01-11 p
\textsc{Meneer Wit}. --- Wel, kom saam met my, my maatjie, ek sal dit vir jou dadelik
verduidelik, want ek is baie lief vir kinders wat hulle wil laat onderrig.

%%K t05-01-12 p
Jy sien daardie werker wat 'n \VUgras{graaf} \textsuperscript{(82)} in die hand hou: hy
is 'n \VUgras{grondwerker} \textsuperscript{(81)}. Hy grawe 'n \VUgras{sloot}
\textsuperscript{(46)} om die waterpype te lê. Hy het ook die grond uitgegrawe op die
plek waar die huis staan, sodat die messelaar die vier \VUgras{mure} kan oprig. Die deel
van daardie mure wat in die grond is, word die \VUgras{fondamente}
\textsuperscript{(2)} genoem. Die ruimte tussen die fondamente vorm die \VUgras{kelder}
\textsuperscript{(3)}. Daardie kelder het 'n klein venster wat die \VUgras{keldervenster}
\textsuperscript{(5)} genoem word, 'n \VUgras{deur} \textsuperscript{(4)} en 'n trappie.

%%K t05-01-13 p
Die \VUgras{messelaar} \textsuperscript{(108)} het aangehou om die mure te bou en
openinge vir die \VUgras{deure} \textsuperscript{(8)} en die \VUgras{vensters}
\textsuperscript{(17)} gelaat.

%%K t05-01-14 p
\textsc{Jan}. --- Maar daardie messelaar sien ek nie!

%%K t05-01-15 p
Mnr. \textsc{Wit}. --- Hy het nou klaar in die huis gewerk. Hy is naby die
\VUgras{stal} \textsuperscript{(54)}, op sy \VUgras{steier} \textsuperscript{(110)}, hy
bou die muur van die \VUgras{waskamer} \textsuperscript{(56)}.

%%K t05-01-16 p
Hy het 'n troffel, 'n bak en 'n \VUgras{skietlood} \textsuperscript{(109)}. Maar jy sal
daardie name nie onthou nie.

%%K t05-01-17 p
\textsc{Jan}. --- Sekerlik wel, want ek sal dadelik by my oom 'n klein troffel en 'n bak
vra, en ek sal self 'n skietlood met 'n toutjie maak.

%%K t05-tit-3 sub
\VUsaut{3.0mm}
\VUpk{10.2pt}{40}{Die materiaal.}
\VUsaut{3.0mm}

%%K t05-01-18 p
\textsc{Die oom}. --- Ha! baie goed. Maar sê my, Jan, wat het 'n mens nodig om 'n huis
te bou?

%%K t05-01-19 p
\textsc{Jan}. --- 'n Mens het \VUgras{kapklippe} \textsuperscript{(59)} en
\VUgras{messelklei} \textsuperscript{(65)} nodig.

%%K t05-01-20 p
\textsc{Die oom}. --- Juis, sien daardie man met die groot hoed, op sy \VUgras{wa}
\textsuperscript{(136)}. Hy is die \VUgras{wadrywer} \textsuperscript{(135)}. Elke dag
bring hy \VUgras{klipblokke} \textsuperscript{(60)} hierheen. Hy laai sy wa op die werf
af, en die \VUgras{klipkappers} \textsuperscript{(83)} kap die blokke en saag hulle met
hulle \VUgras{klipsaag} \textsuperscript{(85)}. 'n \VUgras{Handlanger}
\textsuperscript{(146)} dra hulle na die messelaar wat hulle inpas.

%%K t05-01-21 p
Intussen berei die \VUgras{aanmaker} \textsuperscript{(87)} die messelklei voor. Hy het
\VUgras{sand} \textsuperscript{(62)}, \VUgras{kalk} \textsuperscript{(63)} en
\VUgras{water} \textsuperscript{(64)} nodig. Die kalk is naby hom in 'n \VUgras{sak}
\textsuperscript{(71)}, 'n \VUgras{messelaarskneg} \textsuperscript{(111)} bring vir hom
die sand op 'n \VUgras{kruiwa} \textsuperscript{(112)} en die \VUgras{leerling}
\textsuperscript{(113)} staan by die pomp (58). Hy maak 'n \VUgras{emmer}
\textsuperscript{(114)} vol. Die messelklei sal gou gereed wees. Die
\VUgras{kleidraer} \textsuperscript{(107)} neem dit en dra dit met die leer op tot op die
steier, waar 'n messelaar werk wat daardie kant van die huis klaarmaak.

%%K t05-01-22 p
En nou, hoe word die werker genoem wat aan die \VUgras{dak} \textsuperscript{(31)}
werk?

%%K t05-01-23 p
\textsc{Jan}. --- Hy is die \VUgras{dakdekker} \textsuperscript{(118)}.

%%K t05-tit-4 sub
\VUsaut{3.0mm}
\VUpk{10.2pt}{40}{Die dak van die huis.}
\VUsaut{3.0mm}

%%K t05-01-24 p
Mnr. \textsc{Wit}. --- Ja, maar eers kom die \VUgras{timmerman}
\textsuperscript{(115)}.

%%K t05-01-25 p
Die timmerman maak die \VUgras{dakstoel} \textsuperscript{(35)}. Hy voeg die
\VUgras{balke} \textsuperscript{(37)} met \VUgras{penne} \textsuperscript{(117)}
saam. Daardie werkers, daar bo, met hulle wye fluweelbroeke, is timmermans. Die een hou
'n balkie vas, en die ander sny 'n pen met sy \VUgras{byl} \textsuperscript{(116)}. Die
een wat naby die \VUgras{skoorsteen} \textsuperscript{(41)} is, is die dakdekker. Hy
spyker latte op die (*) leggers, en \VUgras{leiklip} \textsuperscript{(66)} op die
latte. Soms lê hy teëls.

%%K t05-noto-1 noto
\VUnotes{143.2mm}{%
(*) \VUgras{Balk-o} = D. \textit{Balken}, E. \textit{joist}, F. \textit{solive}, I.
\textit{travicello}, R. \textit{balka}, S. Port. \textit{viga}. 'n Tegniese wortel wat
tot nou toe nie aangeneem is nie, maar wat voorgestel behoort te word om die sin van die
Franse woord \textit{solive} weer te gee, wat nie 'n balk F. \textit{poutre} is nie, en
ook nie 'n balkie nie. Dit is 'n vierkantig bewerkte houtstok wat 'n vloer en 'n plafon
dra.%
}

%%K t05-01-26 p
Die dak van die stal is met \VUgras{teëls gedek} \textsuperscript{(55)}, dié van die
huis is met \VUgras{leiklip gedek} \textsuperscript{(32)} en dié van die veranda is van
\VUgras{sink} \textsuperscript{(24)}.

%%K t05-01-27 p
\textsc{Jan}. --- Werk die dakdekker ook aan die sinkdakke?

%%K t05-01-28 p
Mnr. \textsc{Wit}. --- Nee, maar die \VUgras{sinkwerker} \textsuperscript{(121)} of die
\VUgras{loodgieter}, die een wat 'n \VUgras{dakvenster} \textsuperscript{(38)} op die
dak van die huis sit. Hy sit ook die \VUgras{afvoerpype} \textsuperscript{(43)} en die
\VUgras{dakgeute} \textsuperscript{(42)}. Hy soldeer die sink en die blik aanmekaar. Hy
het die \VUgras{weerligafleier} \textsuperscript{(40)} en die \VUgras{weerhaan}
\textsuperscript{(39)} op die \VUgras{gewel} \textsuperscript{(36)} vasgemaak.

%%K t05-01-29 p
\textsc{Jan}. --- Is die huis dan klaar?

%%K t05-tit-5 sub
\VUsaut{3.0mm}
\VUpk{10.2pt}{40}{Die gereedskap.}
\VUsaut{3.0mm}

%%K t05-01-30 p
Mnr. \textsc{Wit}. --- Nie heeltemal nie. Die \VUgras{gipswerker}
\textsuperscript{(99)} bou met \VUgras{bakstene} \textsuperscript{(61)} die mure wat dit
in verskillende kamers verdeel. Hy smeer \VUgras{gips} \textsuperscript{(70)} aan die
\VUgras{plafonne} \textsuperscript{(26)} en die mure. Nou werk hy
aan die plafon van die \VUgras{veranda} \textsuperscript{(33)}. Hy het, soos die
messelaar, 'n \VUgras{troffel} \textsuperscript{(100)} en 'n \VUgras{bak}
\textsuperscript{(101)}.

%%K t05-01-31 p
Daarna maak die skrynwerker die deure, die vensters, die \VUgras{parkette}
\textsuperscript{(16)} en die \VUgras{hortjies} \textsuperscript{(19)}.

%%K t05-01-32 p
Nou werk hy op die werf by sy \VUgras{werkbank} \textsuperscript{(90)}. Hy het 'n
\VUgras{saag} \textsuperscript{(94)} om die hout (69) te saag, 'n \VUgras{skaaf}
\textsuperscript{(91)}, 'n \VUgras{klamp} \textsuperscript{(92)}, 'n \VUgras{beitel}
\textsuperscript{(94 bis)}, 'n \VUgras{passer} \textsuperscript{(95)} en 'n houthamer
\VUgras{(95 bis)}.

%%K t05-01-33 p
\textsc{Jan}. --- Ha! maar dit is lekkerder om te skrynwerk as om te messel. Oom, jy sal
vir my 'n werkbank, 'n saag en 'n skaaf koop, want ek sal binnekort 'n hok vir Medor
maak.

%%K t05-01-34 p
\textsc{Die oom}. --- Ja, seuntjie.

%%K t05-01-35 p
\textsc{Jan}. --- Hoe tevrede is ek! En wie is dit wat naby die voordeur staan?

%%K t05-tit-6 sub
\VUsaut{3.0mm}
\VUpk{10.2pt}{40}{Die verfwerk.}
\VUsaut{3.0mm}

%%K t05-01-36 p
Mnr. \textsc{Wit}. --- Hy is die \VUgras{verwer} \textsuperscript{(102)}. Hy staan op sy
leer met 'n pot \VUgras{verf} \textsuperscript{(103)} en sy \VUgras{kwas}
\textsuperscript{(104)}. Hy verf die hout van die voordeur sodat dit langer hou.

%%K t05-01-37 p
Hy plak ook die papiere in die vertrekke.

%%K t05-01-38 p
Die \VUgras{behanger} \textsuperscript{(107)} wat op 'n \VUgras{leer}
\textsuperscript{(106)} is, op die \VUgras{balkon} \textsuperscript{(28)}, sit 'n
\VUgras{rolgordyn} \textsuperscript{(29)} op.

%%K t05-01-39 p
Die werker wat naby die groot venster staan, is die \VUgras{glasmaker}
\textsuperscript{(96)}. Hy maak die \VUgras{ruite} \textsuperscript{(18)} met
\VUgras{stopverf} \textsuperscript{(73)} vas. Daardie ander werker, wat naby die
kelderdeur kniel, is 'n \VUgras{slotmaker} \textsuperscript{(97)}. Hy sit 'n
\VUgras{slot} \textsuperscript{(6)} in. Sy \VUgras{vyl} \textsuperscript{(98)} lê naby
hom.

%%K t05-01-40 p
\textsc{Jan}. --- Is daardie werker wat met sy \VUgras{hamer} \textsuperscript{(84)} op
ysterwerk slaan ook 'n slotmaker?

%%K t05-01-41 p
Mnr. \textsc{Wit}. --- Om juister te wees, hy is 'n smid (125). Hy smee
\VUgras{ysterwerk} \textsuperscript{(67)} vir die \VUgras{elektrisiën}
\textsuperscript{(124)} wat op die dak van die \VUgras{rytuigskuur}
\textsuperscript{(51)} is. Hy het 'n \VUgras{smidsvuur} \textsuperscript{(126)}, 'n
\VUgras{aambeeld} \textsuperscript{(127)} en 'n \VUgras{tang} \textsuperscript{(128)}.

%%K t05-01-42 p
Die elektrisiën maak die \VUgras{koper} \textsuperscript{(44)} \VUgras{draad} van die
telefoon vas.

%%K t05-tit-7 sub
\VUpk{10.2pt}{40}{Die tuinbou.}
\VUsaut{3.0mm}

%%K t05-01-43 p
\textsc{Die oom}. --- Agterin die \VUgras{werf} \textsuperscript{(45)}, naby die
\VUgras{hek} \textsuperscript{(47)} en die \VUgras{poort} \textsuperscript{(48)}, sien
jy die \VUgras{tuinier} \textsuperscript{(129)}. Hy maak die \VUgras{tuintjie}
\textsuperscript{(49)} gereed. Hy het die grond met sy \VUgras{spitgraaf}
\textsuperscript{(131)} omgespit, hy saai sade en hark met die \VUgras{hark}
\textsuperscript{(130)}. Daarna sal hy met sy \VUgras{gieter} \textsuperscript{(132)}
die \VUgras{beddings} \textsuperscript{(150)} natgooi en die plante sal groei.

%%K t05-01-44 p
Die \VUgras{staljong} \textsuperscript{(133)} wat pas die perd versorg en vir hom kos
gegee het, maak die \VUgras{rytuig} \textsuperscript{(52)} skoon met 'n spons, 'n
\VUgras{handpomp} \textsuperscript{(134)} en 'n emmer water. Daarna sal hy dit weer in
die rytuigskuur inbring en die deur met die dik \VUgras{grendel} \textsuperscript{(53)}
toemaak.

%%K t05-01-45 p
Nou is jy tevrede, dink ek, jy het genoeg verduidelikings gehad.

%%K t05-01-46 p
\textsc{Jan}. --- Ja, oom, maar ek sou graag die binnekant van die huis wou sien.

%%K t05-01-47 p
\textsc{Die oom}. --- Dit sal môre gebeur. Meneer Rooi sal die plan vir jou verduidelik
en jou die naam van elke kamer aandui.

%%K t05-tit-8 sub
\VUsaut{3.0mm}
\VUpk{10.2pt}{40}{Die binne-indeling van die huis.}
\VUsaut{2.4mm}

%%K t05-apar-2 apar
{\centering\textit{(Sien die plan.)}\par}
\VUsaut{2.4mm}

%%K t05-01-48 p
\textsc{Die argitek}. --- Kom, Jan, ek sal jou dadelik die verskillende dele van die
huis laat verken.

%%K t05-01-49 p
Kom ons gaan die \VUgras{grondverdieping} \textsuperscript{(7)} binne. Hier is die
knoppie van die \VUgras{klokkie} \textsuperscript{(11)}. Ons klim die drie
\VUgras{trappies} \textsuperscript{(14)} van die \VUgras{stoep} \textsuperscript{(9)}
op, ons gaan oor die \VUgras{drumpel} \textsuperscript{(10)} van die
\VUgras{voordeur} \textsuperscript{(8)} en hier staan ons aan die voet van die
\VUgras{trap} \textsuperscript{(13)}.

%%K t05-01-50 p
\textsc{Jan}. --- Dit is die gang.

%%K t05-01-51 p
\textsc{Die argitek}. --- Nee, dit is die \VUgras{voorportaal} \textsuperscript{(12)}.
Die \VUgras{gang} \textsuperscript{(a)} is aan die punt. Regs is die \VUgras{sitkamer}
\textsuperscript{(b)} en die \VUgras{eetkamer} \textsuperscript{(c)}, teenoor die
eetkamer is die \VUgras{kombuis} \textsuperscript{(d)} en die \VUgras{bediendekamer}
\textsuperscript{(e)}, dan vooraan die \VUgras{studeerkamer} \textsuperscript{(f)}. Die
\VUgras{veranda} \textsuperscript{(23)} met sy \VUgras{glasdeur} \textsuperscript{(25)}
is naby die sitkamer.

%%K t05-01-52 p
Ons sal nou die trap opklim. Hou jou aan die \VUgras{leuning} \textsuperscript{(15)}
vas en tel die trappies. Daar is veertien. Hier is die \VUgras{gang}
\textsuperscript{(m)}, ons is op die eerste \VUgras{verdieping} \textsuperscript{(27)}.

%%K t05-tit-9 sub
\VUsaut{3.0mm}
\VUpk{10.2pt}{40}{Die eerste verdieping.}
\VUsaut{3.0mm}

%%K t05-01-53 p
Teenoor die trap is die \VUgras{kleinhuisie} \textsuperscript{(20)} en die
\VUgras{kleedkamer} \textsuperscript{(j)}. Regs 'n mooi \VUgras{slaapkamer}
\textsuperscript{(g)} met twee vensters op die \VUgras{voorgewel} \textsuperscript{(1)}
van die huis. Links is die \VUgras{badkamer} \textsuperscript{(1)} en die
\VUgras{gastekamer} \textsuperscript{(i)} met 'n groot \VUgras{alkoof}
\textsuperscript{(h)}. Naby die slaapkamer is die \VUgras{kinderkamer}
\textsuperscript{(k)}. Dit lê aan 'n ruim \VUgras{terras} \textsuperscript{(21)}. Onder
hierdie terras is 'n ander kleinhuisie (20) en, teen die muur, hier is die
\VUgras{klok} \textsuperscript{(22)} wat vir die ete lui.

%%K t05-01-54 p
\textsc{Jan}. --- Kom ons gaan na die \VUgras{tweede verdieping}.

%%K t05-01-55 p
\textsc{Die argitek}. --- Daar is nie 'n tweede verdieping nie. Onder die dak is daar
net \VUgras{solderkamers} \textsuperscript{(33)} en die solder (34). Ons het die
verkenning klaargemaak, ons kan weer afgaan.

%%K t05-01-56 p
\textsc{Jan}. --- Dankie, meneer, dit is baie interessant. Ek sal dadelik vir my vader
alles vertel wat ek gesien het.
\VUsaut{4.0mm}
\VUfilet{22mm}
